\documentclass{amsart}

\usepackage{amsmath}
\usepackage{amssymb}
\usepackage{amsthm}
\usepackage{xcolor}
\usepackage[colorlinks=true,linkcolor=blue,citecolor=blue]{hyperref}

\theoremstyle{plain}
\newtheorem{thm}{Theorem}[section]
\newtheorem{lem}[thm]{Lemma}
\newtheorem{prop}[thm]{Proposition}
\newtheorem{cor}[thm]{Corollary}

\theoremstyle{definition}
\newtheorem{defn}[thm]{Definition}

\theoremstyle{remark}
\newtheorem{rem}[thm]{Remark}

% Editorial macros (never define as no-ops; never drop).
\newcommand{\edit}[1]{\textcolor{red}{#1}}
\newcommand{\note}[1]{\textbf{[#1]}}
\newcommand{\todo}[1]{\textbf{[TODO: #1]}}

\DeclareMathOperator{\Oddsum}{S}

\title{The Sum of the First $n$ Odd Numbers}

\author{A. Author}
\address{Department of Example Studies, Example University}
\email{a.author@example.edu}

\author{B. Coauthor}
\address{Institute for Placeholder Mathematics}
\email{b.coauthor@example.org}

\subjclass[2020]{11A25, 97F60}
\keywords{odd numbers, partial sums, mathematical induction, perfect squares}
\date{}

\begin{document}

\begin{abstract}
We prove that the sum of the first $n$ positive odd numbers equals $n^2$ for every positive integer $n$. The argument is an induction comparing two one-step recurrences. The main result of this paper was obtained with the assistance of the Danus system, an automated proof-search and verification system.
\end{abstract}

\maketitle

\section{Introduction}\label{sec:introduction}

The partial sums of the positive odd numbers form one of the oldest examples in elementary number theory, and the identity recorded below is a standard first exercise in mathematical induction \cite{Exm20}. We give a short, self-contained proof that isolates the single ingredient the argument needs: a comparison of two sequences with equal initial value and equal one-step increments.

For $n \ge 1$, write $\Oddsum(n) = 1 + 3 + 5 + \cdots + (2n-1)$ for the sum of the first $n$ positive odd numbers, with $\Oddsum(1) = 1$. Our result is the following.

\begin{thm}\label{thm:main}
For every integer $n \ge 1$,
\[
  \Oddsum(n) = n^2.
\]
\end{thm}

We fix notation once. Throughout, $n$ denotes a positive integer, and $\Oddsum(n)$ is as defined above. We use no conventions beyond the elementary arithmetic of the integers.

\section{Proof of Theorem~\ref{thm:main}}\label{sec:proof}

We first record two one-step recurrences, then compare them by induction.

\begin{lem}\label{lem:odd-recurrence}
For every integer $n \ge 1$, $\Oddsum(n+1) = \Oddsum(n) + (2n+1)$.
\end{lem}

\begin{proof}
By definition $\Oddsum(n) = \sum_{k=1}^{n} (2k-1)$. The $(n+1)$-st positive odd number is $2(n+1)-1 = 2n+1$, so
\[
  \Oddsum(n+1) = \sum_{k=1}^{n+1} (2k-1) = \left( \sum_{k=1}^{n} (2k-1) \right) + (2n+1) = \Oddsum(n) + (2n+1).
\]
This is the telescoping relation between consecutive partial sums \cite{AC24}.
\end{proof}

\begin{lem}\label{lem:square-recurrence}
For every integer $n \ge 1$, $(n+1)^2 = n^2 + (2n+1)$.
\end{lem}

\begin{proof}
Expanding the square, $(n+1)^2 = n^2 + 2n + 1 = n^2 + (2n+1)$.
\end{proof}

We now prove Theorem~\ref{thm:main}.

\begin{proof}[Proof of Theorem~\ref{thm:main}]
We argue by induction on $n$, in the form recalled in \cite[Section~1]{Exm20}.

\noindent\textbf{Step 1.} The base case holds. For $n = 1$ we have $\Oddsum(1) = 1 = 1^2$.

\noindent\textbf{Step 2.} The inductive step holds. Suppose $\Oddsum(n) = n^2$ for some $n \ge 1$. By Lemma~\ref{lem:odd-recurrence}, $\Oddsum(n+1) = \Oddsum(n) + (2n+1)$, and substituting the inductive hypothesis gives $\Oddsum(n+1) = n^2 + (2n+1)$. By Lemma~\ref{lem:square-recurrence}, $n^2 + (2n+1) = (n+1)^2$. Hence $\Oddsum(n+1) = (n+1)^2$.

By induction, $\Oddsum(n) = n^2$ for every $n \ge 1$.
\end{proof}

\begin{rem}
The main result of this paper was obtained with the assistance of the Danus system, an automated proof-search and verification system. Because automated systems have limitations, it is possible that we have missed related references in the literature, and we welcome comments from experts.
\end{rem}

\begin{thebibliography}{99}

\bibitem[AC24]{AC24}
A. Author and B. Coauthor,
\emph{A note on telescoping sums},
J. Example Math. (2024).

\bibitem[Exm20]{Exm20}
C. Example,
\emph{Elementary induction, revisited},
Example Lecture Notes (2020).

\end{thebibliography}

\end{document}
